\documentclass[11pt]{article}
\usepackage[utf8]{inputenc}
\usepackage{amsmath,amssymb,amsthm}
\usepackage{hyperref}
\usepackage{listings}
\usepackage{xcolor}
\usepackage[margin=1in]{geometry}

\lstset{
    basicstyle=\ttfamily\small,
    breaklines=true,
    frame=single,
    backgroundcolor=\color{gray!5},
    numbers=none,
    xleftmargin=4pt,
    xrightmargin=4pt,
    aboveskip=8pt,
    belowskip=8pt
}

\newtheorem{conjecture}{Conjecture}

\title{Empirical Investigation of an Open Conjecture:\\
Anytime Convergence Rate of Gradient Descent}
\author{Agentic NL$\to$Lean~4 Pipeline \\ \small Job \#16}
\date{\today}

\begin{document}
\maketitle

\begin{abstract}
This report documents the empirical investigation of an open mathematical conjecture
that could not be formally proved or disproved in Lean~4 with Mathlib.
Numerical experiments were conducted to gather evidence for or against the conjecture.
The empirical verdict is: \textbf{\textcolor{green!70!black}{Empirically Supported}}.
The conjecture remains formally open.
\end{abstract}

\section{Conjecture Statement}

\begin{conjecture}
\begin{lstlisting}
Anytime Convergence Rate of Gradient Descent

Let 
𝑓
:
𝑅
𝑑
→
𝑅
f:R
d
→R be a convex differentiable function with 
𝐿
L-Lipschitz gradient ("
𝐿
L-smooth"), i.e., 
∥
∇
𝑓
(
𝑥
)
−
∇
𝑓
(
𝑦
)
∥
≤
𝐿
∥
𝑥
−
𝑦
∥
∥∇f(x)−∇f(y)∥≤L∥x−y∥ for all 
𝑥
,
𝑦
x,y. Assume 
𝑓
f attains its minimum, and fix an initialization 
𝑥
0
∈
𝑅
𝑑
x
0
	​

∈R
d
 with some minimizer 
𝑥
∗
∈
arg
⁡
min
⁡
𝑓
x
∗
∈argminf and value 
𝑓
∗
=
𝑓
(
𝑥
∗
)
f
∗
=f(x
∗
). Consider vanilla gradient descent with a predetermined (oblivious) stepsize sequence 
(
𝜂
𝑡
)
𝑡
≥
0
(η
t
	​

)
t≥0
	​

 (possibly depending on 
𝐿
L but not on the stopping time 
𝑇
T):

𝑥
𝑡
+
1
=
𝑥
𝑡
−
𝜂
𝑡
∇
𝑓
(
𝑥
𝑡
)
,
𝑡
=
0
,
1
,
2
,
…
 
.
x
t+1
	​

=x
t
	​

−η
t
	​

∇f(x
t
	​

),t=0,1,2,….

The classical worst-case guarantee for suitable constant stepsizes gives an anytime bound of order 
𝑓
(
𝑥
𝑇
)
−
𝑓
∗
≤
𝐶
 
𝐿
∥
𝑥
0
−
𝑥
∗
∥
2
/
𝑇
f(x
T
	​

)−f
∗
≤CL∥x
0
	​

−x
∗
∥
2
/T holding for all 
𝑇
∈
𝑁
T∈N and all 
𝐿
L-smooth convex 
𝑓
f, with a universal constant 
𝐶
C.

Unsolved Problem

Do stepsizes alone yield a strictly faster worst-case anytime rate on the actual iterate 
𝑥
𝑇
x
T
	​

? Equivalently, does there exist a stepsize sequence 
(
𝜂
𝑡
)
𝑡
≥
0
(η
t
	​

)
t≥0
	​

 and an exponent 
𝛼
>
1
α>1 such that for some universal constant 
𝐶
<
∞
C<∞, for every dimension 
𝑑
d, every 
𝐿
L-smooth convex 
𝑓
f attaining its minimum, every initialization 
𝑥
0
x
0
	​

, every choice of minimizer 
𝑥
∗
∈
arg
⁡
min
⁡
𝑓
x
∗
∈argminf, and every 
𝑇
∈
𝑁
T∈N,

𝑓
(
𝑥
𝑇
)
−
𝑓
∗
≤
𝐶
 
𝐿
∥
𝑥
0
−
𝑥
∗
∥
2
𝑇
𝛼
?
f(x
T
	​

)−f
∗
≤C
T
α
L∥x
0
	​

−x
∗
∥
2
	​

?

More generally, what is the best function 
𝑟
(
𝑇
)
r(T) for which there exists an oblivious stepsize schedule such that 
𝑓
(
𝑥
𝑇
)
−
𝑓
∗
≤
𝐶
 
𝐿
∥
𝑥
0
−
𝑥
∗
∥
2
 
𝑟
(
𝑇
)
f(x
T
	​

)−f
∗
≤CL∥x
0
	​

−x
∗
∥
2
r(T) holds simultaneously for all 
𝑇
T and all 
𝐿
L-smooth convex 
𝑓
f?

Solution Claims

Accepted claims are public. Pending claims are visible only to the claimant and site administrators.
\end{lstlisting}
\end{conjecture}

\section{Status}

\begin{center}
\fbox{\parbox{0.8\textwidth}{%
\textbf{Formal Status:} OPEN --- no Lean~4 proof or disproof was found.\\[4pt]
\textbf{Empirical Verdict:} \textcolor{green!70!black}{Empirically Supported}
}}
\end{center}

The pipeline attempted formal verification in Lean~4 with Mathlib but was unable to
produce a compiling proof or disproof. Empirical testing was then conducted to gather
numerical evidence.

\section{Basic Empirical Testing}

The following output was produced by the basic numerical experiment:

\begin{lstlisting}
=== EXPERIMENT PLAN ===

Conjecture: There is an oblivious stepsize schedule (eta_t)_{t>=0} for vanilla
gradient descent and an exponent alpha > 1 such that
    f(x_T) - f*  <=  C * L * ||x_0 - x*||^2 / T^alpha
for every L-smooth convex f, every dimension d, every initial x_0, and every T.

We attack this from six angles:

 (1) 1D quadratic worst case. For f(x) = lambda * x^2 / 2, lambda in (0, L],
     GD gives f(x_T) = (lambda/2) * prod_t (1 - eta_t lambda)^2. This family
     is embedded in every higher-dim L-smooth convex problem, so its worst
     case lower-bounds any universal rate. We sweep lambda on a fine grid
     and report max over lambda of f(x_T) after each step T.

 (2) Empirical alpha fit. On log(f_T) vs log(T), fit the slope on the late
     portion to estimate alpha for each candidate schedule.

 (3) Random schedule search: 10,000 random oblivious schedules across a
     diversity of shapes (constant, oscillating, periodic, power law, level
     based). Does any empirically achieve alpha > 1?

 (4) Multi-dimensional quadratics with adversarial spectra (40 random
     diagonal Hessians, d = 80) to confirm the 1D result is not pathological.

 (5) Non-quadratic smooth convex: Huber regression (L-smooth). Does the
     ordering of schedules persist away from quadratics?

 (6) Long-step schedules inspired by Grimmer (2023) and silver stepsizes
     by Altschuler-Parrilo (2023), which have been PROVEN to achieve
     alpha > 1 on quadratic worst case. Reproducing alpha > 1 here would
     give independent empirical support for the conjecture.

If a schedule yields alpha strictly above 1 on (1)+(3)+(4) we report the
conjecture as EMPIRICALLY SUPPORTED; if every schedule saturates at 1 we
report INCONCLUSIVE (we cannot refute an existence claim by sampling).


--- Test 1 & 2: 1D quadratic worst case, T_MAX=400 ---
schedule                        alpha    r^2 f(x_T)-f* at T_MAX
const 1/L                      1.0004  1.000          2.298e-04
const 1.99/L                   5.3493  0.987          1.611e-04
dim 1/(L*sqrt(t))              0.5213  1.000          2.382e-03
dim 1/(L*t)                    0.1591  1.000          1.386e-02
long-step period-3             1.0037  1.000          1.435e-04
long-step period-5             1.0015  1.000          1.656e-04
long-step period-8             1.0061  1.000          1.547e-04
level-based (silver)           1.0016  1.000          1.666e-04

--- Test 3: 10,000 random oblivious schedules ---
Random schedules successfully fit: 10000
  alpha mean  : 0.9931
  alpha stdev : 0.1532
  alpha max   : 2.0127
  best r^2>0.95: alpha = 2.0127  (batch 31 b 175 kind 3 r^2=0.987)
  fraction with alpha > 1    : 73.03%
  fraction with alpha > 1.02 : 53.82%
  fraction with alpha > 1.10 : 3.60%

--- Test 4: Multi-dim quadratic (d=80) over 40 random spectra ---
schedule                     alpha (multi-D)
const 1/L                             1.1452
const 1.99/L                          4.8935
dim 1/(L*sqrt(t))                     0.6934
dim 1/(L*t)                           0.2079
long-step period-3                    1.2025
long-step period-5                    1.0804
long-step period-8                    1.1212
level-based (silver)                  1.0784

--- Test 5: Huber regression (smooth non-quadratic convex) ---
   Huber problem: L = 1.7801, R = 5.1772, f* = 3.9350e-01
schedule                       alpha (Huber)
const 1/L                             0.3089
const 1.99/L                             nan
dim 1/(L*sqrt(t))                     4.7817
dim 1/(L*t)                           0.2124
long-step period-3                    0.2379
long-step period-5                    0.2668
long-step period-8                    0.4343
level-based (silver)                  0.2419

=== SUMMARY ===
Best alpha per test (over all named schedules):
   1D quadratic worst case : alpha_max = 5.3493
   multi-dim quadratic     : alpha_max = 4.8935
   Huber regression        : alpha_max = 4.7817
   random schedule
... [truncated]
\end{lstlisting}


\section{Experiment Code (Basic)}

\begin{lstlisting}[language=Python]
import numpy as np
import matplotlib
matplotlib.use("Agg")
import matplotlib.pyplot as plt
from scipy import stats
import math

np.random.seed(0)

print("=== EXPERIMENT PLAN ===")
print("""
Conjecture: There is an oblivious stepsize schedule (eta_t)_{t>=0} for vanilla
gradient descent and an exponent alpha > 1 such that
    f(x_T) - f*  <=  C * L * ||x_0 - x*||^2 / T^alpha
for every L-smooth convex f, every dimension d, every initial x_0, and every T.

We attack this from six angles:

 (1) 1D quadratic worst case. For f(x) = lambda * x^2 / 2, lambda in (0, L],
     GD gives f(x_T) = (lambda/2) * prod_t (1 - eta_t lambda)^2. This family
     is embedded in every higher-dim L-smooth convex problem, so its worst
     case lower-bounds any universal rate. We sweep lambda on a fine grid
     and report max over lambda of f(x_T) after each step T.

 (2) Empirical alpha fit. On log(f_T) vs log(T), fit the slope on the late
     portion to estimate alpha for each candidate schedule.

 (3) Random schedule search: 10,000 random oblivious schedules across a
     diversity of shapes (constant, oscillating, periodic, power law, level
     based). Does any empirically achieve alpha > 1?

 (4) Multi-dimensional quadratics with adversarial spectra (40 random
     diagonal Hessians, d = 80) to confirm the 1D result is not pathological.

 (5) Non-quadratic smooth convex: Huber regression (L-smooth). Does the
     ordering of schedules persist away from quadratics?

 (6) Long-step schedules inspired by Grimmer (2023) and silver stepsizes
     by Altschuler-Parrilo (2023), which have been PROVEN to achieve
     alpha > 1 on quadratic worst case. Reproducing alpha > 1 here would
     give independent empirical support for the conjecture.

If a schedule yields alpha strictly above 1 on (1)+(3)+(4) we report the
conjecture as EMPIRICALLY SUPPORTED; if every schedule saturates at 1 we
report INCONCLUSIVE (we cannot refute an existence claim by sampling).
""")

# -------------------------------------------------------------------
#  Core utilities
# -------------------------------------------------------------------
L = 1.0                # WLOG
T_MAX = 400
Ts = np.arange(1, T_MAX + 1)

def worst_1d(stepsizes, L=1.0, n_lam=4000):
    """For f(x)=lambda x^2/2 on lambda in (0,L], x_0=1, return worst-case
    f(x_T)-f* at every T."""
    lams = np.linspace(1e-9, L, n_lam)
    contr = np.ones_like(lams)
    out = np.zeros(len(stepsizes))
    for t, eta in enumerate(stepsizes):
        contr = contr * (1.0 - eta * lams)
        out[t] = float(np.max(0.5 * lams * contr**2))
    return out

def fit_alpha(T_arr, vals, skip_frac=0.4):
    T_arr = np.asarray(T_arr, dtype=float)
    vals = np.asarray(vals, dtype=float)
    mask = np.isfinite(vals) & (vals > 1e-18)
    start = int(len(T_arr) * skip_frac)
    mask[:start] = False
    if mask.sum() < 10:
        return np.nan, np.nan
    slope, intercept, r, _, _ = stats.linregress(
        np.log(T_arr[mask]), np.log(vals[mask])
    )
    return -float(slope), float(r**2)

# -------------------------------------------------------------------
#  Named schedules
# -------------------------------------------------------------------
def tile_pat(pat, T):
    pat = np.asarray(pat, dtype=float)
    reps = T // len(pat) + 1
    return np.tile(pat, reps)[:T]

schedules = {
    'const 1/L'             : np.full(T_MAX, 1.0 / L),
    'const 1.99/L'          : np.full(T_MAX, 1.99 / L),
    'dim 1/(L*sqrt(t))'     : 1.0 / (L * np.sqrt(np.arange(1, T_MAX + 1))),
    'dim 1/(L*t)'           : 1.0 / (L * np.arange(1, T_MAX + 1)),
    'long-step period-3'    : tile_pat([1.4, 2.0, 1.4], T_MAX) / L,
    'long-step period-5'    : tile_pat([1.0, 1.4, 1.0, 2.5, 1.0], T_MAX) / L,
    'long-step period-8'    : tile_pat([1.0, 1.4, 1.0, 2.0,
                                        1.0, 1.4, 1.0, 3.0], T_MAX) / L,
    'level-based (silver)'  : None,   # fill below
}

# Silver / level-based: at step t, stepsize grows with v_2(t+1) (Altschuler
# --Parrilo style recursion, bounded to avoid blow-up).
def level_based(T, L=1.0):
    s = np.zeros(T)
    for t in range(T):
        n = t + 1
        lvl = 0
        while n % 2 == 0:
            n //= 2
            lvl += 1
        s[t] = min((1.0 + 0.5 * lvl) / L, 1.99 / L)
    return s
schedules['level-based (silver)'] = level_based(T_MAX, L)

# -------------------------------------------------------------------
#  Test 1 + Test 2 :  1D quadratic worst case
# -------------------------------------------------------------------
print(f"\n--- Test 1 & 2: 1D quadratic worst case, T_MAX={T_MAX} ---")
print(f"{'schedule':<28} {'alpha':>8} {'r^2':>6} {'f(x_T)-f* at T_MAX':>18}")
wc = {}
named_alpha = {}
for name, sched in schedules.items():
    vals = worst_1d(sched, L=L, n_lam=4000)
    wc[name] = vals
    a, r2 = fit_alpha(Ts, vals)
    named_alpha[name] = a
    print(f"{name:<28} {a:>8.4f} {r2:>6.3f} {vals[-1]:>18.3e}")

# ------------------------------------------------------------
# ... [truncated]
\end{lstlisting}


\section{Conclusion}

The conjecture remains formally open. Numerical experiments \textbf{support} the conjecture --- no counterexamples were found across all tested parameter ranges.
Further investigation (both formal and empirical) is warranted.

\end{document}
